% Chapter 4

\chapter{\uppercase{Implementation}}
\label{ch:impl}
This chapter covers the tools we used, the six main algorithms (each with the Maya
example walked through it) and how the engine plugs into Pi.

\section{\uppercase{Implementation Environment}}
The memory service is written in Rust. Tokio gives us the async runtime for the
background workers, Axum handles the API requests and SQLx talks to PostgreSQL.  PostgreSQL 17 with pgvector is
the main store and handles both full-text and vector search, and Neo4j 5.26 holds the
graph copy used for graph-based retrieval.

The models run locally on Ollama. Our extraction model is a fine-tuned
Qwen3-4B-Instruct-2507, which extracts assertions, combines them into inferences and
writes cited answers. Qwen3-Embedding-0.6B makes the vectors. On the agent side we use the Pi agent
harness to capture messages and inject memory. The dashboard (session controls and
the memory graph) is React + TypeScript with Vite and Bootstrap.

Training used Python, PyTorch, Transformers, PEFT and bitsandbytes, with llama.cpp to
make the GGUF. Docker Compose runs the services, and Caddy and ttyd give browser
access to the app and the Pi console.

\section{\uppercase{Algorithms}}
There are six algorithms. Context selection is in M2; entity resolution and versioning
is in M3; consolidation and invalidation is split between M2 (making inferences) and
M3 (marking them stale); and retrieval, rank fusion with packing, and answer generation
are all in M4. 

\subsection{Priority-Based Context Selection}
\label{sec:alg-context}
The extraction model needs to see the assertions already stored about the same people,
otherwise it can't tell a repeat from a correction, but they may not all fit its input
limit. So assertions mentioned in the new messages are ranked first (recency breaks ties) and
a binary search finds the largest prefix that fits.

\begin{algorithm}[tbp]
\caption{Priority-based context selection}
\label{alg:context}
\begin{pseudo}
Input: existing assertions, new messages, input space

IF all existing assertions fit:
    RETURN all assertions
Score assertions by subject and value mentions in new messages
Rank stronger matches first; use recency to break ties
Use binary search to find the largest ranked subset that
    fits
RETURN selected assertions in their original order
\end{pseudo}
\end{algorithm}

For Maya, when E3 (``I moved from Chennai to Bengaluru today\ldots'') comes in, she
has four stored assertions: Chennai, NovaTech, software engineer and remote. Assume for
the example that only three fit. The home city assertion matches E3 on both subject and value,
so it ranks first and the other three follow by recency. Binary search settles on three
assertions, the home city one included, and these go to the extractor with E3.

\subsection{Entity Resolution and Memory Versioning}
This decides which entity a claim is about and what it does to the stored assertions. The
important rule is that a conflicting value never overwrites silently, and the old assertion is
superseded, not deleted, which is what lets us answer the ``previous city'' part of E5.

\begin{algorithm}[tbp]
\caption{Entity resolution and memory versioning}
\label{alg:versioning}
\begin{pseudo}
Input: a new claim and stored memories

Match entity by type and normalised name, or a unique
    stored alias; otherwise create an entity
IF the same assertion exists:
    Attach the new source quote
ELSE IF no conflicting single-value assertion exists:
    Store the claim as a new active assertion
ELSE IF explicit correction AND confidence >= 0.8:
    Supersede the old assertion; activate the new one
ELSE:
    Mark both values as disputed
\end{pseudo}
\end{algorithm}

The E3 claim is Maya, home city Bengaluru, single-valued,
\code{correction=true}, confidence 0.95, quote ``Bengaluru is now my home city.'' Maya is
found, so no new entity. The same assertion doesn't exist (R1 says Chennai), there is a
conflicting single-valued assertion, and the claim is a correction with 0.95 $\geq$ 0.8. So R1
Chennai is superseded with validity 1--10 Oct, and R2 Bengaluru is active from 10 Oct
with E3 as its source.

\subsection{Memory Consolidation and Invalidation}
The model proposes an inference but the code decides whether to save it. Supports
are rechecked just before saving since they may have changed while the model was
running. Invalidation goes the other way and marks dependent inferences stale.

\begin{algorithm}[tbp]
\caption{Memory consolidation and invalidation}
\label{alg:consolidation}
\begin{pseudo}
Input: active assertions about the affected entities

Ask the Extractor SLM to combine related assertions into
    inferences
Require at least two distinct supporting assertions
Check support IDs and the inference's subject
Recheck that supporting assertions are active before saving
Store each inference with its support links

WHEN a supporting assertion changes or is retracted:
    Recursively find dependent inferences
    Mark them as stale
\end{pseudo}
\end{algorithm}

For Maya the inputs are A1 (NovaTech), A2 (software engineer) and A3 (remote). The
model proposes ``Maya works remotely as a software engineer at NovaTech.'' with all
three as supports. It has three distinct supports, the IDs exist, the subject is Maya, and
all three are still active, so it is saved as I1. If A3 were retracted later, I1 would be
found as a dependent and marked stale.

\subsection{Hybrid Retrieval and Graph Traversal}
This returns separate ranked candidate lists. The graph expansion is limited to two
rounds so it doesn't pull in the whole graph. Merging is left to the next algorithm.

\begin{algorithm}[tbp]
\caption{Hybrid retrieval and graph traversal}
\label{alg:retrieval}
\begin{pseudo}
Input: a question and optional date restrictions

Resolve supported date phrases when planning is enabled
Find candidates using keyword and vector searches
Include time-based candidates when a date window exists
Start graph traversal from matching assertions and entities
Expand through entity/support links for two bounded rounds
Filter results by validity and any supplied relevance
    cutoff
RETURN candidate lists for ranking
\end{pseudo}
\end{algorithm}

E5 has no date phrase, so there is no time-based list. Keyword and vector search
find R2 (Bengaluru) and the work assertions. The graph step starts from these and the Maya
entity: \code{MENTIONS} reaches Maya, and the \code{supersedes} link of R2 reaches R1
(Chennai). The graph hits are rechecked in PostgreSQL for status, time and distance, and
the lexical, semantic and graph lists are passed on.

\subsection{Reciprocal Rank Fusion and Context Packing}
Scores from different searches can't be compared, so only ranks are used. A memory
that several lists rank highly beats one that a single list ranks first, and ties go by ID
so the order is deterministic.

\begin{algorithm}[tbp]
\caption{Reciprocal Rank Fusion and context packing}
\label{alg:rrf}
\begin{pseudo}
Input: ranked candidate lists and a context budget

FOR each memory:
    Score = SUM[1 / (60 + rank)] across search lists
Sort by decreasing score; break ties by memory ID
FOR each ranked memory:
    Build a block with the assertion and its source quotes
    Add it if it fits; otherwise consider the next memory
    Stop when the requested number of memories is reached
RETURN the packed context
\end{pseudo}
\end{algorithm}

To illustrate, say R2 is ranked first and R1 second in the lexical, semantic and graph
lists. Then R2 scores $3 \times 1/61 \approx 0.0492$ and R1 scores $3 \times 1/62 \approx 0.0484$, so R2 comes
first and no tie-break is needed. With Pi's settings (\code{top\_k} 10, 2,048 tokens) both
blocks fit easily, and R2 goes out with the E3 quote attached.

\subsection{Source-Grounded Answer Generation}
This decides whether the model's answer is accepted. Bad output is rejected rather
than patched up, so every citation the user sees points to a memory that was actually
retrieved.

\begin{algorithm}[tbp]
\caption{Source-grounded answer generation}
\label{alg:answer}
\begin{pseudo}
Input: a question and retrieved memories

IF no memories were retrieved:
    RETURN "Insufficient evidence."
Ask the extractor SLM to answer from the retrieved memories
Restrict citation IDs to the retrieved memories
Check the response format and citation IDs
Require citations when sufficient evidence is claimed
RETURN the answer, or reject invalid output
\end{pseudo}
\end{algorithm}

For E5 the retrieved memories are R2, R1 and the work assertions, so the model is called.  Only R1, R2 and the work assertion IDs are
allowed as citations; the answer cites R2 for Bengaluru and R1 for Chennai, so it passes.

\section{\uppercase{Integration with the Pi Agent}}
Pi's memory client saves every message, tagged with a namespace and session. When Maya asks E5, Pi sends a recall request for \code{user:maya}
with \code{include\_raw} true, \code{max\_distance} 0.45, \code{top\_k} 10 and \code{max\_tokens} 2,048.  The result is injected only if non-empty and from the same namespace. 
